%% kapa: exam=1 round=19 year=2008 period=2 form=A subject=economics subject_ko=경제원론 seq=12 number=12
%% source: Q-Net 감정평가사 기출문제 (cst00309) artlSeq=1001962 file=extracted/1008/1001962/2008년 제19회 감정평가사 1차시험 기출문제/2교시.hwp
%% conversion: kordoc 4.24.51 -> tools/gen1.py (auto); review=proofread-sonnet; flags=none
\begin{question}{12}
독점적 경쟁시장에 참여하는 기업이 장기균형에서 $Q^{ * }$개의 상품을 생산하고 있다고 하자. $Q^{ * }$에 대한 설명으로 옳은 것은? (단, 생산요소시장은 완전경쟁적이다.)
\begin{choices}
  \item $Q^{ * }$는 가격과 한계비용이 같은 생산량 수준 보다 작다.
  \item $Q^{ * }$는 한계수입과 한계비용이 같은 생산량 수준 보다 크다.
  \item $Q^{ * }$는 한계비용과 평균총비용이 같은 생산량 수준 보다 크다.
  \item $Q^{ * }$는 수요곡선과 평균총비용이 접하는 생산량 수준 보다 작다.
  \item $Q^{ * }$는 평균총비용이 최소화되는 생산량 수준 보다 크다.
\end{choices}
\end{question}
\answer{1}
